\chapter{Xây dựng đồ thị tri thức}
\label{ch:graph}
\section{Quy trình xây dựng}
Hình~\ref{fig:pipeline} mô tả luồng xử lý từ bản ghi chính thức đến đồ thị đã
kiểm tra. Đối sánh công trình tạo bản ghi bài báo được bổ sung siêu dữ liệu.
Nhánh định danh/liên kết tổ chức và nhánh phân loại chủ đề chạy song song:
thiếu bằng chứng về một người không xóa bài chính thức, và nhãn chủ đề không
giải quyết định danh tác giả. Bước tạo đồ thị kết hợp các bản ghi được chấp
nhận, giữ bằng chứng và dẫn xuất quan hệ hợp tác cùng hồ sơ chủ đề.

\begin{figure}[!htbp]\centering
\begin{tikzpicture}[font=\small,
 box/.style={draw=UETBlue,fill=PaleBlue,rounded corners=2pt,align=center,text width=6.3cm,minimum height=1.1cm,inner sep=4pt},
 wide/.style={box,text width=14cm}, arrow/.style={-{Latex[length=2mm]},semithick,color=UETBlue}]
\node[wide] (source) at (0,0) {\textbf{\lang{1. Source collection}{1. Thu thập nguồn}}\\\lang{Official AAAI/IJCAI rosters + cached OpenAlex metadata}{Danh mục AAAI/IJCAI chính thức + siêu dữ liệu OpenAlex lưu đệm}};
\node[wide] (match) at (0,-1.6) {\textbf{\lang{2. Paper matching and cleaning}{2. Đối sánh và làm sạch bản ghi bài báo}}\\\lang{Roster-derived DOI; title and byline-count checks}{DOI từ danh mục; kiểm tra tiêu đề và số tác giả}};
\draw[arrow] (source)--(match);
\node[box,minimum height=1.65cm] (identity) at (-3.8,-3.6) {\textbf{\lang{3--4. People and organizations}{3--4. Tác giả và tổ chức}}\\\lang{Compatible author IDs; ORCID conflicts; paper-specific affiliation triples}{ID tác giả phù hợp; xung đột ORCID; bộ ba liên kết tổ chức theo từng bài}};
\node[box,minimum height=1.65cm] (nlp) at (3.8,-3.6) {\textbf{\lang{5. NLP topic classification}{5. Phân loại chủ đề bằng NLP}}\\\lang{Title/abstract MiniLM embeddings; frozen taxonomy; labels or abstention}{Vector MiniLM từ tiêu đề/tóm tắt; hệ chủ đề cố định; gán nhãn hoặc từ chối}};
\draw[arrow] (match.south -| identity.north)--(identity.north);
\draw[arrow] (match.south -| nlp.north)--(nlp.north);
\node[wide] (graph) at (0,-5.6) {\textbf{\lang{6. Graph materialization}{6. Tạo các quan hệ đồ thị}}\\\lang{Five entity types; observed, inferred, curated and derived edges; evidence links}{Năm loại thực thể; cạnh quan sát, suy ra, biên soạn và dẫn xuất; liên kết bằng chứng}};
\draw[arrow] (identity.south)--(identity.south |- graph.north);
\draw[arrow] (nlp.south)--(nlp.south |- graph.north);
\node[wide] (validate) at (0,-7.25) {\textbf{\lang{7. Validation and storage}{7. Kiểm tra và lưu trữ}}\\\lang{Schema, endpoints, evidence, support counts; Neo4j; analysis and Explorer}{Lược đồ, đầu mút, bằng chứng, số bài hỗ trợ; Neo4j; phân tích và Explorer}};
\draw[arrow] (graph)--(validate);
\end{tikzpicture}
\caption{\lang{Data-to-graph workflow. Identity/affiliation handling and topic assignment are parallel dependencies. Frozen v4 retains the seven historical DOI exceptions described in Chapter~\ref{ch:data}.}{Quy trình từ dữ liệu đến đồ thị. Xử lý định danh/liên kết tổ chức và gán chủ đề là hai nhánh song song. V4 cố định vẫn giữ bảy trường hợp DOI lịch sử nêu ở Chương~\ref{ch:data}.}}
\label{fig:pipeline}\end{figure}


Đồ thị được lưu dưới dạng JSON và nạp vào cơ sở dữ liệu
Neo4j Community 5.26 riêng biệt \citep{neo4jDocs}. Các phép chiếu
được định nghĩa riêng phục vụ phân tích ngoại tuyến; Explorer cho phép đọc
thực thể, quan hệ và bằng chứng nguồn. Mọi số liệu trong báo cáo đều được tính
từ cùng một đồ thị và cấu hình phân tích.

\section{Phân loại chủ đề bằng NLP}
Dự án dùng xử lý ngôn ngữ tự nhiên (NLP) để nối bài báo với chủ đề nghiên cứu.
Với mỗi công trình đối sánh, đầu vào là tiêu đề nối với tóm tắt được khôi phục;
hệ thống hỗ trợ chỉ dùng tiêu đề khi thiếu tóm tắt. Cùng một mô hình cố định
\modelid{sentence-transformers/all-MiniLM-L6-v2} mã hóa văn bản bài báo và mô tả
chủ đề lá thành vector \citep{reimers2019sentencebert,minilmModelCard}. Mô tả
kết hợp tên, định nghĩa, các khái niệm bao hàm và từ đồng nghĩa của chủ đề.
Mã phiên bản mô hình được ghim trong cấu hình kho mã nguồn.

Bộ phân loại so sánh các vector chuẩn hóa bằng độ tương đồng cosin. Nhãn tốt
nhất được nhận khi điểm đạt ngưỡng hiệu chỉnh trên tập phát triển
$\tau\approx0.2709$. Có thể nhận thêm tối đa hai nhãn lá nếu chúng cũng vượt
ngưỡng và cách điểm tốt nhất không quá $\delta=0.02$. Nếu không đủ điều kiện,
bộ phân loại từ chối gán nhãn. Cách này tạo cạnh \texttt{HAS\_TOPIC} đa nhãn
thay vì buộc mọi bài vào một nhóm. Câu hỗ trợ trong tóm tắt được chọn theo độ
tương đồng ngữ nghĩa và lưu cùng bộ định vị nguồn; với đầu vào chỉ có tiêu đề,
tiêu đề là bằng chứng. Điểm tương đồng không phải xác suất đã hiệu chuẩn.

Hệ chủ đề có bốn chủ đề cha và 16 chủ đề lá. Các cạnh \texttt{SUBTOPIC\_OF}
được biên soạn riêng với nhãn bài báo do mô hình suy ra. Hồ sơ chủ đề của nhà
nghiên cứu và tổ chức được dẫn xuất sau đó từ các bài có bằng chứng được chấp
nhận. Mô hình và ngưỡng giữ nguyên; báo cáo này không chạy lại suy luận chủ đề.

\section{Phân giải thực thể và ánh xạ liên kết tổ chức}
\textbf{Bài báo.} ID danh mục chính thức xác định bản ghi kỷ yếu; DOI, tiêu đề
và danh sách tác giả nối nó với công trình OpenAlex. Ứng viên mơ hồ không được
chấp nhận. Lưu ý DOI lịch sử được trình bày ở Chương~\ref{ch:data}.

\textbf{Nhà nghiên cứu.} Tại từng vị trí tác giả, tên chính thức phù hợp với
bản ghi tác giả OpenAlex cung cấp ID nhà nghiên cứu. Khóa nội bộ chứa ID tác
giả OpenAlex; tên chuẩn hóa trùng nhau không đủ để hợp nhất người. ORCID được
giữ khi có, nhưng không bắt buộc. Nhiều ORCID có mã kiểm tra hợp lệ cùng gắn
với một ID nhà cung cấp gây giữ lại ID đó trên toàn tập. Hơn một vị trí không
khớp tên hoặc ID tác giả lặp trong bài có thể khiến toàn bộ ánh xạ tác giả của
bài bị giữ lại. Vì vậy, đối sánh công trình thành công không có nghĩa là đã
phân giải mọi tác giả trên bài.

\textbf{Tổ chức.} Liên kết tổ chức (affiliation) được xét theo từng bài. Tổ chức phải có ID OpenAlex và xuất hiện ở đúng bản ghi tác
giả được chấp nhận. Bản ghi tổ chức giữ khu vực hoạt động và quốc gia đã chuẩn
hóa. Quy trình lưu từng quan sát nhà nghiên cứu--bài báo--tổ chức trước khi tổng
hợp cạnh \texttt{AFFILIATED\_WITH}. Chuỗi đơn vị công tác gốc là bằng chứng
nguồn; hồ sơ hiện tại hoặc danh sách tổ chức của cả bài không được sao chép cho
mọi tác giả. Quy tắc bảo vệ giữ lại các tập nhiều tổ chức giống hệt nhau xuất
hiện đáng ngờ ở nhiều tác giả. Tổ chức hoặc khu vực chưa xác định được giữ ở
trạng thái chưa rõ.

\section{Lược đồ đồ thị và ngữ nghĩa quan hệ}
Năm loại nút là nhà nghiên cứu, bài báo, tổ chức, chủ đề và hội nghị.
Bảng~\ref{tab:schema} tóm tắt mười quan hệ. Bốn loại bằng chứng phân biệt siêu
dữ liệu quan sát, suy luận mô hình, hệ từ vựng biên soạn và quan hệ dẫn xuất từ
cạnh cơ sở được chấp nhận.

\begin{table}[!htbp]\centering\small
\caption{\lang{The ten relationships and their support rules. Arrows specify direction.}{Mười quan hệ và quy tắc hỗ trợ. Mũi tên chỉ chiều quan hệ.}}
\label{tab:schema}
\setlength{\tabcolsep}{4pt}\renewcommand{\arraystretch}{1.12}
\begin{tabularx}{\textwidth}{>{\raggedright\arraybackslash}p{4.1cm}>{\raggedright\arraybackslash}p{3.7cm}Y}\toprule
\lang{Predicate / class}{Quan hệ / loại} & \lang{Endpoints}{Hai đầu quan hệ} & \lang{Meaning}{Ý nghĩa}\\\midrule
\texttt{AUTHORED} & \lang{Researcher $\to$ Paper}{Nhà nghiên cứu $\to$ Bài báo} & \lang{Observed: accepted byline identity.}{Quan sát: định danh tác giả được chấp nhận.}\\
\texttt{AFFILIATED\_WITH} & \lang{Researcher $\to$ Organization}{Nhà nghiên cứu $\to$ Tổ chức} & \lang{Observed: aggregate paper-specific affiliations.}{Quan sát: tổng hợp liên kết tổ chức theo từng bài.}\\
\texttt{PUBLISHED\_IN} & \lang{Paper $\to$ Venue}{Bài báo $\to$ Hội nghị} & \lang{Observed: official proceedings entry.}{Quan sát: bản ghi kỷ yếu chính thức.}\\
\texttt{CITES} & \lang{Paper $\to$ Paper}{Bài báo $\to$ Bài báo} & \lang{Observed: in-scope citing-to-cited reference.}{Quan sát: trích dẫn tới bài trong phạm vi.}\\
\texttt{HAS\_TOPIC} & \lang{Paper $\to$ Topic}{Bài báo $\to$ Chủ đề} & \lang{Inferred: one to three leaf labels, or abstention.}{Suy ra từ mô hình: một đến ba nhãn lá, hoặc không gán.}\\
\texttt{SUBTOPIC\_OF} & \lang{Topic $\to$ Topic}{Chủ đề $\to$ Chủ đề} & \lang{Curated: leaf-to-parent taxonomy.}{Biên soạn: chủ đề lá thuộc chủ đề cha.}\\
\texttt{COAUTHORED\_WITH} & \lang{Researcher $\leftrightarrow$ Researcher}{Nhà nghiên cứu $\leftrightarrow$ Nhà nghiên cứu} & \lang{Derived: shared accepted core paper.}{Dẫn xuất: cùng bài thuộc tập lõi, tác giả được chấp nhận.}\\
\texttt{COLLABORATED\_WITH} & \lang{Organization $\leftrightarrow$ Organization}{Tổ chức $\leftrightarrow$ Tổ chức} & \lang{Derived: distinct authors with aligned affiliations.}{Dẫn xuất: tác giả khác nhau có liên kết tổ chức tương ứng.}\\
\texttt{RESEARCHES\_TOPIC} & \lang{Researcher $\to$ Topic}{Nhà nghiên cứu $\to$ Chủ đề} & \lang{Derived: at least two distinct core papers.}{Dẫn xuất: ít nhất hai bài phân biệt thuộc tập lõi.}\\
\makecell[l]{\texttt{CONTRIBUTES\_TO\_}\\\texttt{TOPIC}} & \lang{Organization $\to$ Topic}{Tổ chức $\to$ Chủ đề} & \lang{Derived: at least three distinct core papers.}{Dẫn xuất: ít nhất ba bài phân biệt thuộc tập lõi.}\\\bottomrule
\end{tabularx}\end{table}

\FloatBarrier

Mỗi cạnh ngữ nghĩa có ID xác định, trạng thái chấp nhận, ID hai đầu, tham chiếu
bằng chứng và thông tin về phương pháp xây dựng. Cặp hợp tác đối xứng chỉ lưu một lần
theo thứ tự ID chuẩn tắc và được phân tích như quan hệ vô hướng. Xác nhận từ
nhiều nguồn hoặc bản sao đảo chiều không tạo cạnh mới.

Với bài có $k$ tác giả chính thức, mỗi cặp đồng tác giả được chấp nhận nhận
trọng số phân bổ $1/(k-1)$. Tác giả chưa phân giải vẫn nằm trong $k$. Hợp tác
tổ chức cần hai tác giả khác nhau hỗ trợ hai tổ chức khác nhau trên cùng bài;
một người có hai đơn vị công tác không tự tạo cạnh hợp tác tổ chức. Nếu bài
hỗ trợ $q$ cặp tổ chức đủ điều kiện, mỗi cặp nhận $1/q$. Cả hai quan hệ đều giữ
số bài hỗ trợ phân biệt. \texttt{RESEARCHES\_TOPIC} cần hai bài phân biệt thuộc
tập lõi; \texttt{CONTRIBUTES\_TO\_TOPIC} cần ba bài. Bằng chứng lặp cho một bài
không làm tăng số đếm.

\Needspace{18cm}
\section{Ví dụ xây dựng đồ thị từ bài báo}
Xét \emph{ECG ODE-GAN: Learning Ordinary Differential Equations of ECG Dynamics
via Generative Adversarial Learning} (AAAI 2021).
\href{https://ojs.aaai.org/index.php/AAAI/article/view/16086}{Bản ghi chính thức}
trở thành \sourcepath{paper:aaai:2021:16086}. DOI
\href{https://doi.org/10.1609/aaai.v35i1.16086}{\texttt{10.1609/aaai.v35i1.16086}}
đối sánh với \href{https://openalex.org/W3176135934}{OpenAlex W3176135934}, với
tiêu đề khớp hoàn toàn và cùng ba tác giả.

Các quan hệ tác giả được chấp nhận nối bài báo với Tomer Golany (A5022512211, vị trí 1) và Kira Radinsky
(A5029708595, vị trí 3). Daniel Freedman, tác giả ở giữa trên danh sách chính
thức, không có cạnh tác giả được chấp nhận trong đồ thị này. Hai quan sát theo
bài nối Golany và Radinsky với Technion (I174306211). Bộ phân loại gán chủ đề
\emph{Mô hình sinh}, với điểm cosin lưu lại là 0.3753 và đoạn tóm tắt về học
động lực học thông qua mạng sinh đối kháng.

\begin{figure}[!htbp]\centering
\begin{tikzpicture}[font=\small,
 box/.style={draw=UETBlue,fill=PaleBlue,rounded corners=2pt,align=center,text width=4.0cm,minimum height=0.9cm,inner sep=4pt},
 arrow/.style={-{Latex[length=1.8mm]},color=UETBlue},lab/.style={font=\footnotesize,fill=white,inner sep=1.5pt}]
\node[box,text width=5.5cm] (org) at (0,3.4) {Technion\\\texttt{I174306211}};
\node[box] (a) at (-4.4,1.6) {Tomer Golany\\\texttt{A5022512211}};
\node[box] (b) at (4.4,1.6) {Kira Radinsky\\\texttt{A5029708595}};
\node[box,text width=5.5cm] (paper) at (0,-0.2) {ECG ODE-GAN (AAAI 2021)\\\texttt{paper:aaai:2021:16086}};
\node[box] (venue) at (-4.4,-2.2) {AAAI\\\texttt{venue:aaai}};
\node[box] (topic) at (4.4,-2.2) {\lang{Generative modeling}{Mô hình sinh}\\{\footnotesize\texttt{topic:}\\[-2pt]\texttt{generative-modeling}}};
\draw[arrow] (a.north)--node[lab,above]{R2}(org.west);
\draw[arrow] (b.north)--node[lab,above]{R2}(org.east);
\draw[arrow] (a.south)--node[lab,below]{R1}(paper.west);
\draw[arrow] (b.south)--node[lab,above]{R1}(paper.east);
\draw[dashed,color=UETGold,thick] (a.east)--node[lab,above]{R7}(b.west);
\draw[arrow] (paper.south west)--node[lab,below]{R3}(venue.north);
\draw[arrow,dashed] (paper.south east)--node[lab,above]{R5}(topic.north);
\end{tikzpicture}
\par\smallskip
{\footnotesize\lang{Relationships}{Quan hệ}: R1 \texttt{AUTHORED}; R2 \texttt{AFFILIATED\_WITH};
R3 \texttt{PUBLISHED\_IN};\\
R5 \texttt{HAS\_TOPIC}; R7 \texttt{COAUTHORED\_WITH}.\par}
\caption{\lang{Accepted relationships for the example. IDs are shortened to provider IDs where shown. Solid arrows are observed; dashed topic arrow is inferred; the dashed undirected line is derived. Affiliation edges aggregate paper-specific observations; both displayed observations exist for this paper.}{Các quan hệ được chấp nhận trong ví dụ. Một số ID được rút gọn thành ID nhà cung cấp. Mũi tên liền là quan sát; mũi tên nét đứt tới chủ đề là suy ra; đường vô hướng nét đứt là dẫn xuất. Cạnh liên kết tổ chức tổng hợp các quan sát theo bài; cả hai quan sát được thể hiện đều tồn tại ở bài này.}}
\label{fig:example}\end{figure}

\FloatBarrier

Các cạnh tác giả, đơn vị công tác, hội nghị, chủ đề và đồng tác giả trong hình
đều tồn tại trong đồ thị. Bài này đóng góp $1/(3-1)=0.5$ vào trọng số phân bổ của cặp
đồng tác giả được chấp nhận; cạnh có thể tổng hợp thêm bài khác. Hai tác giả
đều liên kết với một tổ chức, nên ví dụ không tạo hợp tác giữa hai tổ chức.
Một bài riêng lẻ cũng chưa đủ lập cạnh hồ sơ chủ đề cho nhà nghiên cứu hoặc tổ
chức. ID đầy đủ, ID cạnh được chấp nhận và hai quan sát liên kết tổ chức nằm tại
\sourcepath{report/latex/data/construction-example.json}.
